\textbf{Время:} 02:00

\textbf{Тезис:} средний квантовый вес модельных пар отличается от экспериментального отношения к mixed-event фону.

\textbf{Визуал:} отобранная пара → фаза → вес → среднее в импульсном бине.

\textbf{Текст:} Для каждой пары берём импульсы и точки испускания. Из их разностей получаем фазу и назначаем вес: единица плюс косинус. Складываем веса в бине и делим на число тех же пар. Результат — средний вес.

Здесь нет двух независимых счётчиков. В числитель и знаменатель входят одни и те же пары, поэтому это другая конструкция, чем экспериментальное смешивание событий. Ошибку отношения независимых пуассоновских гистограмм использовать нельзя.

Нужны точки последнего взаимодействия или рождения конечного пиона при распаде. Общая поздняя поверхность вывода содержит уже свободный пролёт; её координаты не равны точкам испускания. Формат данных генератора необходимо проверить.

Косинус учитывает свободные плоские волны. Для кулоновского и сильного взаимодействий нужен другой вес с соответствующей волновой функцией. Начать проверку полезно с искусственного гауссова источника известной ширины, затем переходить к транспортным событиям. Неопределённость итоговых параметров проверяют с учётом структуры событий.

\textbf{Переход:} Посмотрим на конкретный результат моего освоения этой обработки.
